% ---------------------------------------------------------------------
\section{核靶}\label{sec:disall-nuclear}

原子核是色单态束缚态，DIS 的证明不用到靶的其他性质；核的特殊之处是线度随质量数 $A$ 增长，使某些幂次修正被增强，并在小 $x$ 处带来相干的多重散射。

\begin{proposition}[核靶上的 DIS \tagA]\label{prop:disall-nuclear}
设靶为质量数 $A$ 的原子核。
\begin{enumerate}[label=(\alph*),leftmargin=2em,itemsep=1pt,topsep=2pt]
\item 在领头幂次，核靶上的 DIS 与强子靶上的 DIS 具有相同的因子化（\cref{thm:disall-factorization}）：以核 PDF 代替强子 PDF，系数函数不变。
\item （核增强的扭度 4 修正，部分，\tagP）在高扭度的推广因子化中\cite{QiuSterman1991a}，核内的多重散射给出的扭度 4 修正被核的线度增强，
相对大小 $\propto A^{1/3}\Lambda^2/Q^2$\cite{Luo:1993ui,Luo:1994np}；据\cite{Luo:1993ui}的摘要，由光生双喷注数据估计的扭度 4 核部分子分布约为 $0.05$--$0.1\,\mathrm{GeV}^2$ 乘以典型的扭度 2 分布。
\item 对强子—核碰撞（例如 pA 中的 Drell--Yan），初态相互作用可能使因子化形式失效，除非束流能量大于一个随靶长度增长的尺度（据\cite{Bodwin:1988fs}的摘要；较早的讨论见\cite{BodwinBrodskyLepage1981}）；
在只有一个强子、末态被求和的 DIS 中，这一初态机制不出现。
\item 小 $x$ 处的核遮蔽可以在领头扭度中，借助 Gribov--Glauber 理论与衍射 DIS 的因子化来描述\cite{Frankfurt:2011cs}。
\item （小 $x$ 的饱和区，猜想，\tagC）在足够小的 $x$ 与大核中，饱和效应可能使领头扭度的描述失效（\cref{ch:smallx}）。
\end{enumerate}
非微扰地位：\cref{asm:intro-bridge}（MA1、MA2）。
\uses{thm:disall-factorization,def:pdf-quark,def:pdf-gluon,asm:intro-bridge}
% errtype: E-S
\end{proposition}

\begin{proofsketch}
(a) 推导只用到靶是色单态的物理态；核 PDF 由\cref{def:pdf-quark,def:pdf-gluon}的同一组算符在核态中定义（通常按每个核子归一化），强子层次同样依赖于\cref{asm:intro-bridge}。
(b) 推广的高扭度因子化中\cite{QiuSterman1991a,QiuSterman1991b}，领头 $1/Q^2$ 修正是硬部分与扭度 2、扭度 4 矩阵元的卷积。直观上，扭度 4 矩阵元含有一个沿光锥方向、不带振荡相因子的位置积分（两对场之间的距离）；
两对场属于不同核子时，这一积分延伸到整个核的直径 $\propto A^{1/3}$，这就是核增强的来源；比例系数只能由数据估计\cite{Luo:1993ui}。(c) DIS 中没有入射强子旁观者的初态相互作用，末态相互作用按\cref{thm:disall-factorization}的论证被求和。
(d) 遮蔽理论结合了 Gribov--Glauber 理论、QCD 因子化与 HERA 衍射 DIS 的分析\cite{Frankfurt:2011cs}（衍射 DIS 见\cref{ch:break}）。(e) 饱和区见\cref{ch:smallx}。
\end{proofsketch}

\paragraph{地位与缺口}
(a) 继承\cref{thm:disall-factorization}的全部瓶颈（\cref{asm:appN-collins,asm:appN-ward-product,asm:appN-cutsum}等），全局核 PDF 拟合都以领头扭度因子化为前提\cite{Eskola:2021nhw,Kovarik:2015cma,AbdulKhalek:2022fyi}。
真正未被控制的是幂次修正对核尺寸的依赖：(b) 只是模型估计（[P]），推广的高扭度因子化在核中没有全阶证明；(e) 是猜想（[C]），它与共线因子化的衔接没有被证明。
